% Apsis of Collision

The apsis of collision is the distance of closest approach during a collision.
The reduced apsis of collision $R_0$ is the solution of \cref{eq:sn_apsis}.
For actually solving this equation with Newtons method, it is convenient to define the function $f(R) = R \, g^2(R)$ so that $R_0$ is the solution of:
%
\begin{equation}
    \label{eq:sn_apsis_eq}
    f(R) = R - \frac{1}{\varepsilon} \Phi(R) - \frac{P^2}{R} = 0.
\end{equation}
%
For screening functions given by a sum of exponentials, \cref{eq:sn_sum_exp}, and also for the finite-range screening function \cref{eq:sn_NLHlin}, the first derivative is always negative, \(\Phi'(R)<0\), and the second derivative is always positive, \(\Phi''(R)>0\).
It follows that \(f'(R)>0\) and \(f''(R)<0\) for all values of \(R\).
As a consequence, Newton’s method is guaranteed to converge if the initial condition \(R_{0,0}\) is chosen to the left of the solution \(R_0\) (\(R_{0,0} < R_0\)).

\subsubsection{Obtaining an initial condition}
%%%%%%%%%%%%%%%%%%%%%%%%%%%%%%%%%%%%%%%%%%%%%%
%
For purely repulsive potentials \(R_0 > P\). Therefore, \(P\) is to the left of \(R_0\), and
%
\begin{equation}
    \label{eq:sn_apsis_ic_p}
    R_0^\mathrm{(0,P)} = P
\end{equation}
%
is a possible initial condition, guaranteeing convergence of Newton’s method.
It turns out that convergence with initial condition \cref{eq:sn_apsis_ic_p} is slow for large scattering angles (i.e., small impact parameters $P$).
Therefore, a second lower boundary to $R_0$ is desirable.

For a head-on collision ($P = 0$), \cref{eq:sn_apsis_eq} reduces to
%
\begin{equation}
    \label{eq:sn_apsis_ic_eps}
    R_0^{(0,\varepsilon)} - \frac{ 1 }{ \varepsilon } \Phi(R_0^{(0,\varepsilon)}) = 0 ~.
\end{equation}
%
where we have replaced \(R\) by \(R_0^{(0,\varepsilon)}\), since we want to use this equation to obtain an alternative initial condition for $R_0$.
This equation also can be solved by iteration only.
However, it is easy to produce a table \(\varepsilon = \varepsilon(R_0^{(0,\varepsilon)})\) by calculating $\varepsilon$ from given values of $R_0^{(0,\varepsilon)}$, and to obtain a value of \(R_0^{(0,\varepsilon)}\) for a given $\varepsilon$ by interpolation in the inverse table \(R_0^{(0,\varepsilon)}(\varepsilon)\).
If the interpolation is done cubic, the derivatives \(\mathrm{d} R_0^{(0,\varepsilon)} / \mathrm{d} \varepsilon\) need to be tabulated as well. They are obtained by differentiating \cref{eq:sn_apsis_ic_eps} with respect to $\varepsilon$ and reinserting $\varepsilon$ from \cref{eq:sn_apsis_ic_eps}:
%
\begin{equation}
    \frac{ \mathrm{d} R_0^{(0,\varepsilon)} }{ \mathrm{d} \varepsilon } =
        - \frac{ (R_0^{(0,\varepsilon)})^2 }{ \Phi(R_0^{(0,\varepsilon)}) - R_0^{(0,\varepsilon)} \, \Phi'(R_0^{(0,\varepsilon)}) } ~.
\end{equation}
%
Note that $\Phi'(R) < 0$ and therefore $\mathrm{d} R_0^{(0,\varepsilon)} / \mathrm{d} \varepsilon < 0$ for all $R_0^{(0,\varepsilon)}$.%
%
\footnote{It can be shown that $\mathrm{d}^2 R_0^{(0,\varepsilon)} / \mathrm{d} \varepsilon^2 > 0$. Therefore it is advisable to use cubic rather than linear interpolation, since linear interpolation would in general overestimate $R_0^{(0,\varepsilon)}$, which could result in $R_0^{(0,\varepsilon)} > R_0$.
}

To set up the table, it seems more convenient to select values for $\varepsilon$ than for $R_0^{(0,\varepsilon)}$, since the range of \(\varepsilon\) required is known to be between the implant and the trajectory cutoff energy.
It is also possible to choose the \(\varepsilon\) limits generously, e.g. by \(\varepsilon_\mathrm{min} = 10^{-8}\) and \(\varepsilon_\mathrm{max} = 10^4\), which then would allow to use a precomputed table for all simulations.

If we desire to determine $R_0^{(0,\varepsilon)}(\varepsilon_i)$ for $\varepsilon_i = \varepsilon_\mathrm{max}/2^i$, $i = 0 \ldots n$, and $\varepsilon_\mathrm{max}/2^n \le \varepsilon_\mathrm{min} < \varepsilon_\mathrm{max}/2^{n-1}$, we can start with determining the \(R_0^{(0,\varepsilon)}\) value corresponding to \(\varepsilon_{max}\) by Newton iteration of \cref{eq:sn_apsis_eq} with $P=0$.
An initial condition for the iteration, which is to the left of the solution, is \(R_0^{(0,\varepsilon,0)} = 1/\varepsilon_\mathrm{max}\), which follows from \cref{eq:sn_apsis_ic_eps} together with \(\Phi(R_0^{(0,\varepsilon,0)}) \le 1\). Once \(R_0^{(0,\varepsilon)}(\varepsilon_\mathrm{max})\) has been determined, \(R_0^{(0,\varepsilon)}(\varepsilon_\mathrm{max}/2)\) is determined in the same way, using \(R_0^{(0,\varepsilon)}(\varepsilon_\mathrm{max})\) as initial condition, then \(R_0^{(0,\varepsilon)}(\varepsilon_\mathrm{max}/4)\) using \(R_0^{(0,\varepsilon)}(\varepsilon_\mathrm{max}/2)\) as initial condition and so on, until \(\varepsilon_\mathrm{max}/2^n \le \varepsilon_\mathrm{min}\).

$R_0^{(0,\varepsilon)}$ can also be taken as an initial condition for $P > 0$. This can be seen by
bringing the second and third term of \cref{eq:sn_apsis_eq} to the right-hand side and writing $R_0$ instead of $R$ (since the equation is fulfilled for $R = R_0$):
%
\begin{equation*}
    R_0 = \frac{1}{\varepsilon} \Phi(R_0) + \frac{P^2}{R_0} .
\end{equation*}
%
Similarly, we can rewrite \cref{eq:sn_apsis_ic_eps}
%
\begin{equation*}
    R_0^{(0,\varepsilon)} = \frac{1}{\varepsilon} \Phi(R_0^{(0,\varepsilon)}) .
\end{equation*}
%
We can easily prove that $R_0^{(0,\varepsilon)} \ge R_0$ is not possible for any $P > 0$. If $R_0^{(0,\varepsilon)} \ge R_0$ were true for some $P > 0$, because of $\mathrm{d} \Phi / \mathrm{d} R < 0$
%
\begin{equation}
    R_0 = \frac{1}{\varepsilon} \Phi(R_0) + \frac{P^2}{R_0}
        \ge \frac{1}{\varepsilon} \Phi(R_0^{(0,\varepsilon)}) + \frac{P^2}{R_0}
        = R_0^{(0,\varepsilon)} +  \frac{P^2}{R_0}
\end{equation}
%
for that $P$.
Since $P^2 / R_0 > 0$, this would mean that $R_0 > R_0^{(0,\varepsilon)}$ which contradicts the assumption $R_0^{(0,\varepsilon)} \ge R_0$.
Consequently, $R_0^{(0,\varepsilon)}$ must be smaller than $R_0$ for all $P > 0$.
So when $R_0^{(0,\varepsilon)}$ is used as an initial condition for the Newton iteration, the iteration is guaranteed to converge.

Now, as we can use both $R_0^{(0,P)}$ and $R_0^{(0,\varepsilon)}$ as initial conditions, we can choose the larger one
%
\begin{equation}
    \label{eq:sn_apsis_ic}
    R_0^{(0)} = \max( R_0^{(0,P)}, R_0^{(0,\varepsilon)} ) ,
\end{equation}
%
and we will do so as it is closer to the final solution $R_0$ than the smaller one.
\Cref{eq:sn_apsis_ic} can then be used as the initial condition in a Newton iteration
%
\begin{equation}
    \label{eq:sn_apsis_newton}
    R_0^{(i+1)} = R_0^{(i)} - f(R_0^{(i)})/f'(R_0^{(i)})
\end{equation}
%
It guarantees convergence of the iterations for screening functions with $\Phi'(R) < 0$ and $\Phi''(r) > 0$.

\subsubsection{Convergence}
%%%%%%%%%%%%%%%%%%%%%%%%%%%
%
\Cref{fig:iter-apsis} shows the number of iterations required to reach a relative accuracy of \(10^{-3}\) between the last two calculated values of \(R_0\) for the universal ZBL and the NLHlin potential.
Usually this means that the last value is accurate to \(10^{-6}\), assuming quadratic convergence has been reached.

\begin{figure}[htbp]
    \centering
    \includegraphics[scale=0.41]{ZBLuniv_iter_apsis.pdf} ~~
    \includegraphics[scale=0.41]{NLHlin_iter_apsis.pdf}
    \caption{Number of iterations required to reach a relative accuracy of \(10^{-3}\) in \(R_0\) as a function of energy and impact parameter.
    Left: Universal ZBL potential.
    Right: NLHlin potential for As-Si using \cref{eq:sn_aNLH} as the screening length.
    The reduced impact parameters $P$ have been chosen as to cover the same range of unscaled impact parameters $p$.
    }
    \label{fig:iter-apsis}
\end{figure}

For combination of low energies and small impact parameters as well as for combinations of high energies and large impact parameters only 1-2 iterations are necessary. This is because in these cases either $R_0^{(0,P)}$ or $R_0^{(0,\varepsilon)}$ already is a good approximation to $R_0$. In other combination of $\varepsilon$ and $P$ up to 4 iterations are necessary.
